\section{数据链路层}

\subsection{数据链路层的功能}

\begin{mydef}{数据链路层}{dll-def}
数据链路层在物理层提供的比特流基础上，为网络层提供\textbf{帧}（Frame）的透明传输。主要功能：
\begin{enumerate}
    \item \textbf{成帧（Framing）}：将比特流划分为帧。
    \item \textbf{差错控制}：检测并（可选地）纠正传输错误。
    \item \textbf{流量控制}：协调发送方和接收方的速率（速率匹配）。
    \item \textbf{介质访问控制（MAC）}：广播网络中决定谁可以使用共享信道。
\end{enumerate}
\end{mydef}

\subsection{组帧方法}

\begin{conclusion}{四种组帧方法}{framing-methods}
\begin{itemize}
    \item \textbf{字符计数法}：帧首部包含帧中字符数。缺点：计数字段一旦出错，后续全部帧错位。
    \item \textbf{字符填充法（字节填充）}：用特殊字符标记帧边界，数据中出现的边界字符用转义字符 ESC 替换。
    \item \textbf{零比特填充法（Bit Stuffing）}：帧边界标记为 \texttt{01111110}（6 个连续 1）。发送方每 5 个连续的 1 后插入 0，接收方每 5 个连续的 1 后删除 0。\textbf{HDLC 和 PPP 使用。}
    \item \textbf{违规编码法}：利用物理层的冗余编码（如曼彻斯特编码中非法跳变）。仅适用于物理层编码特殊的网络。
\end{itemize}
\end{conclusion}

\subsection{差错控制}

\begin{formula}{CRC 循环冗余校验}{crc}
\textbf{原理：}将 $k$ 位数据后拼接 $r$ 位校验位，构成 $n=k+r$ 位码字，使其能被生成多项式 $G(x)$ 整除。接收方用同一 $G(x)$ 除码字，余数非零表示出错。

\textbf{408 计算方法：}给定数据比特串 $M$ 和生成多项式 $G(x)$（如 $G(x)=x^3+x+1$，即 1011）：
\begin{enumerate}
    \item $M$ 后补 $r$ 个 0（$r$ 是 $G(x)$ 的次数），得 $2^r M$。
    \item $2^r M$ 除以 $G(x)$（模 2 除法，即异或运算），得 $r$ 位余数 $R$。
    \item 发送码字 = $M\parallel R$（$M$ 后拼接 $R$）。
\end{enumerate}
\end{formula}

\begin{attention}
\textbf{408 CRC 计算题的几个注意点：}
\begin{itemize}
    \item CRC 是\textbf{检错码}（非纠错码）。能检测所有单比特错误和绝大多数突发错误。
    \item 模 2 除法 = 异或（不进位、不借位）。
    \item $r$ 次生成多项式产生的 CRC 校验位为 $r$ 位。
\end{itemize}
\end{attention}

\subsection{介质访问控制（MAC）子层}

\begin{conclusion}{信道划分协议}{channel-partition}
\begin{itemize}
    \item \textbf{频分复用（FDM）}：各用户占用不同频段。
    \item \textbf{时分复用（TDM）}：各用户分配不同时隙。统计 TDM（STDM）按需分配，效率更高。
    \item \textbf{波分复用（WDM）}：光的频分复用。
    \item \textbf{码分复用（CDM/CDMA）}：每用户分配唯一码片序列，各用户同时同频传输。接收方用码片序列内积提取数据。
\end{itemize}
\end{conclusion}

\subsection{CSMA/CD 协议（以太网核心）}

\begin{conclusion}{CSMA/CD 协议}{csma-cd}
CSMA/CD（Carrier Sense Multiple Access with Collision Detection）是以太网的核心协议。

\textbf{工作流程}（四句话）：
\begin{enumerate}
    \item \textbf{先听后发}：发送前先监听信道，信道空闲则发送。
    \item \textbf{边发边听}：发送过程中持续监听。
    \item \textbf{冲突停发}：检测到冲突立即停止发送。
    \item \textbf{随机重发}：等待随机时间后重新尝试（\textbf{二进制指数退避算法}）。
\end{enumerate}

\textbf{争用期（Contention Period）}：$2\tau$（$\tau$ 为端到端传播时延）。在争用期内若未检测到冲突，则后续不会冲突。

\textbf{最短帧长}：$L_{\min} = 2\tau \times \text{数据率}$。若帧过短（小于 $L_{\min}$），发送方在检测到冲突前就已发完，无法重传。

\textbf{二进制指数退避：}第 $k$ 次重传（$k\leqslant 10$）的等待时间为从 $[0, 2^k-1]$ 中随机选一个 $r$，等待 $r\times 2\tau$。$k>10$ 后 $r$ 范围锁定为 $[0,2^{10}-1]$。重传 16 次仍失败则丢弃帧。
\end{conclusion}

\begin{attention}
\textbf{408 CSMA/CD 必考计算：}
\begin{enumerate}
    \item 给定数据率和最大距离，求最短帧长。
    \item 给定最短帧长和数据率，求最大传输距离。
    \item 二进制指数退避的等待时间和 $k$ 值的对应关系。
\end{enumerate}
\end{attention}

\subsection{以太网与交换机}

\begin{conclusion}{以太网帧格式}{ethernet-frame}
\[
\begin{array}{|c|c|c|c|c|c|c|}
\hline
\text{前导码} & \text{目的MAC} & \text{源MAC} & \text{类型} & \text{数据} & \text{填充} & \text{CRC} \\
\text{(8B)} & \text{(6B)} & \text{(6B)} & \text{(2B)} & \text{(46-1500B)} & & \text{(4B)} \\
\hline
\end{array}
\]
最短帧长 64B = 8+6+6+2+46+4（数据至少 46B 凑足最短帧长）。
\end{conclusion}

\begin{conclusion}{交换机与网桥}{switch-bridge}
\begin{itemize}
    \item \textbf{网桥（Bridge）}：工作在数据链路层，根据 MAC 地址转发帧。可隔离冲突域，但无法隔离广播域。
    \item \textbf{交换机（Switch）}：多端口网桥。每个端口是一个独立的冲突域。
    \item \textbf{自学习算法}：交换机通过源 MAC 地址学习端口-地址映射，构建\textbf{交换表（MAC 地址表）}。未知目的地址的帧\textbf{泛洪}（flood）到除接收端口外的所有端口。
\end{itemize}

\begin{tablebox}{冲突域 vs 广播域}
\centering
\begin{tabular}{|l|c|c|}
\hline
设备 & 隔离冲突域？& 隔离广播域？\\
\hline
中继器/集线器（物理层）& 否 & 否 \\
网桥/交换机（数据链路层）& \textbf{是} & 否 \\
路由器（网络层）& 是 & \textbf{是} \\
\hline
\end{tabular}
\end{tablebox}
\end{conclusion}

\subsection{ARP 协议}

\begin{mydef}{ARP}{arp}
ARP（Address Resolution Protocol）将 IP 地址映射为 MAC 地址。工作过程：主机广播 ARP 请求（「谁的 IP 是 xxx？请告诉你的 MAC」），目标主机单播 ARP 响应。结果缓存于 ARP 表（有老化时间）。
\end{mydef}

\subsection{PPP 协议}

\begin{conclusion}{PPP}{ppp}
点对点协议（Point-to-Point Protocol）是目前最广泛使用的数据链路层协议（如电话线拨号、PPPoE）。特点：
\begin{itemize}
    \item 不提供可靠传输（不编号、不确认），差错检测仅用 CRC。
    \item 面向字节，使用字节填充（转义字符 \texttt{0x7D}）。
    \item 支持多种网络层协议（LCP + NCP）。
\end{itemize}
\end{conclusion}

\subsection{408 高频考点速查}

\begin{conclusion}{数据链路层 — 408常考点}{cn3-keypoints}
\begin{enumerate}
    \item \textbf{CSMA/CD}：最短帧长计算、争用期 $2\tau$ 的概念、二进制指数退避。
    \item \textbf{交换机自学习}：给定帧转发序列和初始空交换表，写出交换表项的变化过程。
    \item \textbf{CRC 校验}：给定生成多项式，计算 CRC 校验码（模 2 除法）。
    \item \textbf{冲突域 vs 广播域}：各设备隔离哪个域。
    \item \textbf{ARP}：工作原理、是否跨网络（ARP 仅在同一广播域内有效）。
\end{enumerate}
\end{conclusion}

\newpage
\section{习题}

\subsection{基础过关题}

\begin{enumerate}

\item \textbf{（2009年·选择题·第36题）} 数据链路层采用后退N帧（GBN）协议，发送方已经发送了编号为0~7的帧。当计时器超时时，若发送方只收到0、2、3号帧的确认，则发送方需要重发的帧数是
\begin{center}
\begin{tabular}{ll}
(A) 2 & (B) 3 \\[6pt]
(C) 4 & (D) 5
\end{tabular}
\end{center}

\ifshowsolutions\par\noindent\textbf{解析：} GBN协议使用累计确认。收到3号帧的确认表示0-3号帧均已正确接收，只需重发4-7号帧，共4帧。\boxed{\textbf{答案：}(C)}\fi

\vspace{8pt}

\item \textbf{（2010年·选择题·第35题）} 某局域网采用CSMA/CD协议，数据传输速率为10 Mbps，信号传播速率为$2\times 10^8$ m/s，主机甲和主机乙相距2 km。若主机甲和主机乙发送数据时发生冲突，则从开始发送数据到检测到冲突的最短时间和最长时间分别是
\begin{center}
\begin{tabular}{ll}
(A) 0, $20\ \mu$s & (B) 0, $10\ \mu$s \\[6pt]
(C) $10\ \mu$s, $20\ \mu$s & (D) $10\ \mu$s, $10\ \mu$s
\end{tabular}
\end{center}

\ifshowsolutions\par\noindent\textbf{解析：} $\tau=2\text{ km}/(2\times 10^8\text{ m/s})=10\ \mu$s。最短：双方同时发送（0），最长：一方发到达另一方前另一方才发（$2\tau=20\ \mu$s）。\boxed{\textbf{答案：}(A)}\fi

\vspace{8pt}

\item \textbf{（2011年·选择题·第35题）} 对于100 Mbps的以太网交换机，当输出端口无排队，以直通交换方式转发一个以太网帧时，引入的转发时延至少是
\begin{center}
\begin{tabular}{ll}
(A) 0\ \mu s & (B) 0.48\ \mu s \\[6pt]
(C) 5.12\ \mu s & (D) 121.44\ \mu s
\end{tabular}
\end{center}

\ifshowsolutions\par\noindent\textbf{解析：} 直通交换只接收目的MAC地址（6 B）就开始转发。$6\times 8/10^8=0.48\ \mu$s。\boxed{\textbf{答案：}(B)}\fi

\vspace{8pt}

\item \textbf{（2013年·选择题·第36题）} 下列介质访问控制方法中，可能发生冲突的是
\begin{center}
\begin{tabular}{ll}
(A) CDMA & (B) CSMA/CD \\[6pt]
(C) TDMA & (D) FDMA
\end{tabular}
\end{center}

\ifshowsolutions\par\noindent\textbf{解析：} CSMA/CD是竞争型协议，存在冲突。CDMA/TDMA/FDMA是信道划分协议，无冲突。\boxed{\textbf{答案：}(B)}\fi

\vspace{8pt}

\item \textbf{（2018年·选择题·第35题）} 数据$M=101001$，生成多项式$G(x)=x^3+x^2+1$（即1101），则CRC校验码为
\begin{center}
\begin{tabular}{ll}
(A) 001 & (B) 101 \\[6pt]
(C) 011 & (D) 110
\end{tabular}
\end{center}

\ifshowsolutions\par\noindent\textbf{解析：} $M$后补3个0得$101001000$。模2除$1101$：$101001000\div 1101$，余数$=001$。\boxed{\textbf{答案：}(A)}\fi

\end{enumerate}

\subsection{真题风格与综合训练}

\begin{enumerate}[resume]

\item \textbf{（2021年·选择题·第35题）} 交换机通过自学习算法建立交换表。某交换机初始交换表为空，先后收到发往不同目的MAC地址的帧。以下关于自学习的说法，正确的是
\begin{center}
\begin{tabular}{ll}
(A) 交换机通过目的MAC地址学习 & (B) 交换机通过源MAC地址学习 \\[4pt]
(C) 交换机只在收到广播帧时学习 & (D) 交换表项永不过期
\end{tabular}
\end{center}

\ifshowsolutions\par\noindent\textbf{解析：} 交换机通过\textbf{源MAC地址}和接收端口学习，建立端口-地址映射。表项有老化时间。\boxed{\textbf{答案：}(B)}\fi

\vspace{8pt}

\item \textbf{（真题风格·综合题）} 某局域网通过交换机连接4台主机A、B、C、D，分别连接端口1-4。交换表初始为空。按顺序：A→B, B→A, C→A, D→C。描述每步后交换表变化及帧的转发方式（明确转发、泛洪或过滤）。

\ifshowsolutions\par\noindent\textbf{解析：}
(1) A→B：表空，无B的映射→泛洪到端口2,3,4。学习A(端口1)。
(2) B→A：表中有A在端口1→明确转发到端口1。学习B(端口2)。
(3) C→A：表中有A在端口1→明确转发到端口1。学习C(端口3)。
(4) D→C：表中有C在端口3→明确转发到端口3。学习D(端口4)。\fi

\end{enumerate}

\vspace{8pt}
\noindent\textbf{拓展阅读：}
\begin{itemize}
    \item Kurose \& Ross 第6章——CSMA/CD、交换机自学习的完整讨论。
    \item 谢希仁 第3章——PPP协议和CSMA/CD的国内标准表述。
\end{itemize}
